\section{Introduction and setup}

\subsection{Introduction}
\label{sec:intro}
% Member 1

\paragraph{The problem.}
A classifier is usually judged by its accuracy, the fraction of labels it gets right. Accuracy
quietly assumes that every mistake costs the same, and that the system must always commit to a label.
Neither is true in most applications. In this project we build a binary \emph{decision} system rather
than a pure classifier. For each observation it chooses the action with the smallest expected cost,
and when it is too unsure it may refer the observation to a person instead of guessing.

\paragraph{The application.}
We study a spam filter. Each email is either spam ($y=1$) or a real email ($y=0$), and the filter
chooses one of three actions: deliver it to the inbox, send it to the spam folder, or hold it for
human inspection. The two errors are very different. A spam in the inbox wastes a few seconds; a real
email in the spam folder may never be read, so we charge it ten times as much ($\cfp=10$ against
$\cfn=1$). Inspection always gives the right answer but costs a person's time ($\ci$). The filter
therefore needs more than a label: it needs the probability $\post=\Prob(y=1\mid x)$ that an email is
spam, and a rule that turns this probability into the cheapest action.

\paragraph{Why accuracy is the wrong target.}
The accuracy-based rule sends an email to the spam folder whenever $\post>1/2$. With our costs the
cheapest rule (Section~\ref{sec:decision}) does so only when $\post\ge0.97$, inspects emails with
$0.3<\post<0.97$, and delivers the rest. The two rules disagree on every email in between, so a model can
have the higher accuracy and still lose more real email. Three further effects widen the gap:
\begin{itemize}[itemsep=1pt,topsep=2pt]
  \item The cost-based rule reads $\post$ at thresholds far from $1/2$, so it needs probabilities that
        are \emph{calibrated} there. Accuracy only checks which side of $1/2$ each probability falls on.
  \item The share of spam at deployment can differ from the share in the training data. That shifts every
        posterior, even though spam and real emails themselves look the same as before.
  \item Inspection helps only when it is cheap enough, and how cheap depends on both error costs.
\end{itemize}

\paragraph{Roadmap.}
Section~\ref{sec:decision} fixes the loss matrix, derives the action that minimises the posterior
expected loss, and gives explicit thresholds with and without an inspect option.
Section~\ref{sec:models} presents the two classifiers we compare: a Gaussian generative classifier
(with shared or class-specific covariance) and logistic regression, both written from scratch.
Section~\ref{sec:task2} compares accuracy-based and loss-minimising decisions over a range of cost
ratios, with confusion matrices and average realised loss.
Section~\ref{sec:calibration} measures probability quality with log-loss, the Brier score and
reliability plots, and explains why correct labels alone do not guarantee good decisions.
Section~\ref{sec:prior-shift} changes the class proportions at deployment while keeping $p(x\mid y)$
fixed, derives the posterior-odds correction for known priors, and tests it.
Section~\ref{sec:accuracy-cost} shows a setting where higher accuracy comes with higher cost, traces
realised cost and rejection rate as the inspection cost varies, and determines when allowing rejection
improves the system. Section~\ref{sec:conclusion} summarises the findings.
% TODO (Member 1): once Sections 4-7 have results, add one sentence with the headline numbers.

\subsection{Experimental setup}
\label{sec:setup}
% Member 6 (C1)

\paragraph{Application, actions and losses.}
We filter email. Label $y=1$ means spam and $y=0$ a real email. For each email the filter
chooses one of three actions: deliver to the inbox ($a=0$), send to the spam folder ($a=1$), or
hold for human inspection ($a=r$; coded as \texttt{2}). Correct decisions cost nothing. Sending
a real email to spam costs $\cfp$, delivering spam costs $\cfn$, and inspection costs $\ci$
whatever the label. We assume the inspector always finds the true label and routes the email
correctly, so no further loss follows an inspection. Unless stated otherwise $\cfp=10$ and $\cfn=1$;
all experiments sweep $\ci$ or the ratio $\cfp/\cfn$.

\paragraph{Data.}
\emph{Spambase}~\cite{spambase} has $4{,}601$ emails, of which $1{,}813$ ($39.4\%$) are spam, and
57 non-negative features: 48 word frequencies, 6 character frequencies and 3 run-length
statistics of capital letters. The features are heavily skewed and $77\%$ of all feature values are
exactly zero, which matters for the Gaussian models (Section~\ref{sec:models}).
The \emph{synthetic} data are two-class Gaussian samples with chosen means, covariances and
prior $\pio$. Because $p(x\mid y)$ is known we can compute the exact posterior and Bayes risk,
and change the class proportions exactly (Section~\ref{sec:prior-shift}). Section~\ref{sec:accuracy-cost}
also uses a one-dimensional special case with closed-form answers.

\paragraph{Splits and error bars.}
Each data set is split into stratified $60/20/20$ train, validation and test parts
($2{,}761/921/919$ emails for Spambase, each with the overall spam rate). Models are fitted on the
training part. Every choice (regularisation strength, temperature, anything else) is made on the
validation part. The test part is used only to report results. We repeat everything over 20 splits
(seeds $0$--$19$) and report the mean with a standard-error bar.

\paragraph{Preprocessing.}
Standardisation (zero mean, unit variance), optionally after the transform $\log(1+x)$, uses
statistics of the training part only and is then applied unchanged to validation and test data.

\paragraph{Prior shift.}
To simulate a deployment prior $\pio'$ we resample test emails \emph{within each class}, so
$p(x\mid y)$ is unchanged and only the label proportions move. By default we sample without
replacement and keep the largest set the data allow.

\paragraph{Code and reproducibility.}
All models are written from scratch in Python with NumPy; scikit-learn is used only to cross-check
numbers. The interface above is implemented by three shared modules: \texttt{data.py},
\texttt{decision.py} and \texttt{metrics.py}. A unit test checks the decision rule of
Section~\ref{sec:decision} against a brute-force
minimisation of the posterior expected loss on $10{,}000$ random posteriors and loss matrices,
covering cases both with and without an inspect band. The command
\texttt{python -m project4.run\_all} regenerates every figure and table in this report from scratch.
